% GENERATED FILE. Edit canonical lessons, then run npm run generate:books.
% canonical-source-hash: fnv1a64:af426395005a5fe5
% canonical-lessons: SA-C180-vyadhah, SA-C180-caurah, SA-C180-yacakah, SA-C180-panditah, SA-C180-vidvan

\chapter{Hunter, Thief, Beggar, Scholar, Learned person}
\label{ch:sa-hunter-thief-beggar-scholar-learned-person}
\hlchaptermodality{180}

\begin{chapteropening}
\textbf{By the end of this chapter:} \emph{I can say a hunter, a thief, a beggar, a scholar and a learned person.}

Complete the last lesson of chapter 180: I can say a hunter, a thief, a beggar, a scholar and a learned person.
\end{chapteropening}

\section[\texorpdfstring{\sk{व्याधः} (vyādhaḥ)}{व्याधः (vyādhaḥ)}]{\sk{व्याधः} (vyādhaḥ) — a hunter}

\label{lesson:SA-C180-vyadhah}



\begin{quote}
Before the new one: say the Sanskrit for a lawyer, then the Sanskrit for a writer.
\end{quote}

\subsection*{You'll want to know: \sk{व्याधः}}
\textbf{\sk{व्याधः}} — \emph{vyādhaḥ} — \textquotedblleft{}a hunter\textquotedblright{}.

\begin{grammarlens}[title={What you've built}]
One more word for this chapter.
\end{grammarlens}

\subsection*{Your turn}
\begin{itemize}
\raggedright
  \item \emph{Say it:} \emph{vyādhaḥ}
  \item \emph{Say it:} \emph{vyādhaḥ}, once more
  \item \emph{Say it:} say \emph{vyādhaḥ}
  \item \emph{Recall:} say the Sanskrit for a lawyer, then the Sanskrit for a writer, then say \emph{vyādhaḥ} again
\end{itemize}

\begin{tcolorbox}[breakable,colback=teal!4,colframe=teal!35!black,title={Before you move on}]
What does \sk{व्याधः} mean? (\textquotedblleft{}A hunter\textquotedblright{}.) Say it once more.
\end{tcolorbox}
\section[\texorpdfstring{\sk{चौरः} (cauraḥ)}{चौरः (cauraḥ)}]{\sk{चौरः} (cauraḥ) — a thief}

\label{lesson:SA-C180-caurah}



\begin{quote}
Before the new one: say the Sanskrit for a writer, then the Sanskrit for a hunter.
\end{quote}

\subsection*{You'll want to know: \sk{चौरः}}
\textbf{\sk{चौरः}} — \emph{cauraḥ} — \textquotedblleft{}a thief\textquotedblright{}.

\begin{grammarlens}[title={What you've built}]
One more word for this chapter.
\end{grammarlens}

\subsection*{Your turn}
\begin{itemize}
\raggedright
  \item \emph{Say it:} \emph{cauraḥ}
  \item \emph{Say it:} \emph{cauraḥ}, once more
  \item \emph{Say it:} say \emph{cauraḥ}
  \item \emph{Recall:} say the Sanskrit for a writer, then the Sanskrit for a hunter, then say \emph{cauraḥ} again
\end{itemize}

\begin{tcolorbox}[breakable,colback=teal!4,colframe=teal!35!black,title={Before you move on}]
What does \sk{चौरः} mean? (\textquotedblleft{}A thief\textquotedblright{}.) Say it once more.
\end{tcolorbox}
\section[\texorpdfstring{\sk{याचकः} (yācakaḥ)}{याचकः (yācakaḥ)}]{\sk{याचकः} (yācakaḥ) — a beggar}

\label{lesson:SA-C180-yacakah}



\begin{quote}
Before the new one: say the Sanskrit for a hunter, then the Sanskrit for a thief.
\end{quote}

\subsection*{You'll want to know: \sk{याचकः}}
\textbf{\sk{याचकः}} — \emph{yācakaḥ} — \textquotedblleft{}a beggar\textquotedblright{}.

\begin{grammarlens}[title={What you've built}]
One more word for this chapter.
\end{grammarlens}

\subsection*{Your turn}
\begin{itemize}
\raggedright
  \item \emph{Say it:} \emph{yācakaḥ}
  \item \emph{Say it:} \emph{yācakaḥ}, once more
  \item \emph{Say it:} say \emph{yācakaḥ}
  \item \emph{Recall:} say the Sanskrit for a hunter, then the Sanskrit for a thief, then say \emph{yācakaḥ} again
\end{itemize}

\begin{tcolorbox}[breakable,colback=teal!4,colframe=teal!35!black,title={Before you move on}]
What does \sk{याचकः} mean? (\textquotedblleft{}A beggar\textquotedblright{}.) Say it once more.
\end{tcolorbox}
\section[\texorpdfstring{\sk{पण्डितः} (paṇḍitaḥ)}{पण्डितः (paṇḍitaḥ)}]{\sk{पण्डितः} (paṇḍitaḥ) — a scholar}

\label{lesson:SA-C180-panditah}



\begin{quote}
Before the new one: say the Sanskrit for a thief, then the Sanskrit for a beggar.
\end{quote}

\subsection*{You'll want to know: \sk{पण्डितः}}
\textbf{\sk{पण्डितः}} — \emph{paṇḍitaḥ} — \textquotedblleft{}a scholar\textquotedblright{}.

\begin{grammarlens}[title={What you've built}]
One more word for this chapter.
\end{grammarlens}

\subsection*{Your turn}
\begin{itemize}
\raggedright
  \item \emph{Say it:} \emph{paṇḍitaḥ}
  \item \emph{Say it:} \emph{paṇḍitaḥ}, once more
  \item \emph{Say it:} say \emph{paṇḍitaḥ}
  \item \emph{Recall:} say the Sanskrit for a thief, then the Sanskrit for a beggar, then say \emph{paṇḍitaḥ} again
\end{itemize}

\begin{tcolorbox}[breakable,colback=teal!4,colframe=teal!35!black,title={Before you move on}]
What does \sk{पण्डितः} mean? (\textquotedblleft{}A scholar\textquotedblright{}.) Say it once more.
\end{tcolorbox}
\section[\texorpdfstring{\sk{विद्वान्} (vidvān)}{विद्वान् (vidvān)}]{\sk{विद्वान्} (vidvān) — a learned person}

\label{lesson:SA-C180-vidvan}



\begin{quote}
Before the new one: say the Sanskrit for a beggar, then the Sanskrit for a scholar.
\end{quote}

\subsection*{You'll want to know: \sk{विद्वान्}}
\textbf{\sk{विद्वान्}} — \emph{vidvān} — \textquotedblleft{}a learned person\textquotedblright{}.

\begin{grammarlens}[title={What you've built}]
One more word for this chapter.
\end{grammarlens}

\subsection*{Your turn}
\begin{itemize}
\raggedright
  \item \emph{Say it:} \emph{vidvān}
  \item \emph{Say it:} \emph{vidvān}, once more
  \item \emph{Say it:} say \emph{vidvān}
  \item \emph{Recall:} say the Sanskrit for a beggar, then the Sanskrit for a scholar, then say \emph{vidvān} again
\end{itemize}

\begin{tcolorbox}[breakable,colback=teal!4,colframe=teal!35!black,title={Before you move on}]
What does \sk{विद्वान्} mean? (\textquotedblleft{}A learned person\textquotedblright{}.) Say it once more.
\end{tcolorbox}
